% DRAFT (LLM-drafted from the protocol documents and the claims ledger; the author reviews and confirms).
\section{The equal-evidence protocol}
\label{sec:protocol}

The protocol has five parts. None is new on its own; the point is to apply all of them to every scored run.

\paragraph{1. Identical, hash-pinned requests.} For each comparison we fix a set of \emph{request bodies}: the
instruction, the options with their names or descriptions, the retrieved examples in order, and the test input. Each
body is serialised once, its SHA-256 is written to a manifest, and every system receives that body (translated only
by the minimal adapter its interface requires, which is itself versioned). Two systems are said to be compared on
\emph{equal evidence} only when their manifests match item by item. In the 24-shot setting this separates two
sources of difference that published numbers conflate: the model and the retrieval rule that chose its examples.

\paragraph{2. Pre-registration before every scored run.} Before a run touches test data, a commit fixes the
configuration, the selection procedure (on data disjoint from test), the primary and secondary analyses, the
multiple-comparison families and the failure handling (for example, a failed API call counts as wrong for accuracy
and is excluded, with a count, from calibration). Later deviations are written as dated amendments before the runs
they govern. The repository history makes the order verifiable: for our headline run, the protocol, the selection
and the single test run are separate commits, and the test-time dump is touched by exactly one commit
(Appendix~\ref{app:protocol}).

\paragraph{3. Paired exact tests.} Every accuracy comparison is an exact McNemar test on the same test items,
reported with its discordant counts $b$ (first system right, second wrong) and $c$ (the reverse). Within a declared
family we apply Holm's step-down correction; post hoc families are labelled as such. Calibration is compared with
paired bootstrap intervals over items (10,000 resamples; 2,000 for ECE) % numbers-ok: protocol constants, jev_swap_analysis.py
for ECE with 10 equal-width bins, the Brier score summed over classes, and log loss where the returned probabilities permit
it (\S\ref{sec:find-rounding}).

\paragraph{4. Release code first; a parity gate for any other engine.} Each open model is scored with its release
code when feasible. When throughput requires a different serving engine, the engine must first reproduce the
release code item by item on a full benchmark (argmax agreement and maximum absolute probability difference are
recorded); a model that fails the gate is scored only with its release code. Over-length requests follow each
model's release code, and the policy is recorded per model, because systems that truncate, reject or silently
shorten long requests can otherwise be ranked on different inputs.

\paragraph{5. Replication before comparison.} When a closed system's published result exists with per-item
predictions, we first re-send its exact request bodies to the pinned version and test agreement item by item. Only a
system whose published result we can reproduce is used as a reference.

\paragraph{What the protocol does not control.} It cannot verify what a model was trained on. We record each
model's documented training data per dataset and report zero-shot results for exposed or undocumented combinations
separately (\S\ref{sec:find-exposure}). It also does not equalise compute: models are compared at their released
sizes, and cost is reported separately (\S\ref{sec:cost}).
